% !TEX root = 概率论期中复习.tex
% 自动抽取自 chapters/chapter2.tex；原例题及其解答保留。

\clearpage\section{随机变量及其分布函数}\label{k:cdf}
\label{sec:random-variables-and-distribution-functions}

\begin{mydefinition}{随机变量}{random-variable}
给定概率空间 $(\Omega,\mathscr F,P)$。若映射 $X:\Omega\to\mathbb R$ 满足：对任意 $B\in\mathscr B(\mathbb R)$，都有
\[
X^{-1}(B):=\{\omega\in\Omega:X(\omega)\in B\}\in\mathscr F,
\]
则称 $X$ 为该概率空间上的\textbf{随机变量}。随机变量通常用 $X,Y,Z$ 或 $\xi,\eta$ 等表示。
\end{mydefinition}
\origin{原笔记第2.1节，定义“随机变量”；源码第9行。}

\begin{proposition}{随机变量的半直线判别}{random-variable-halfline-criterion}
映射 $X:\Omega\to\mathbb R$ 是随机变量，当且仅当对任意 $t\in\mathbb R$ 都有
\[
\{\omega\in\Omega:X(\omega)\leq t\}\in\mathscr F.
\]
\end{proposition}
\origin{原笔记第2.1节，命题“随机变量的半直线判别”；源码第19行。}

\subsection{例题：示性函数}\label{note:c2:e1}
\basis{原笔记第2.1节，例“示性函数”；源码\texttt{chapters/chapter2.tex}第26行。}
设 $A\in\mathscr F$，定义
\[
I_A(\omega)=
\begin{cases}
1,&\omega\in A,\\
0,&\omega\notin A.
\end{cases}
\]
称 $I_A$ 为事件 $A$ 的\textbf{示性函数}。由于
\[
\{I_A\leq t\}=
\begin{cases}
\varnothing,&t<0,\\
A^{\mathrm c},&0\leq t<1,\\
\Omega,&t\geq1,
\end{cases}
\]
上述集合均属于 $\mathscr F$，所以 $I_A$ 是随机变量。并且
\[
P(I_A=1)=P(A),\qquad P(I_A=0)=P(A^{\mathrm c}).
\]

\begin{mydefinition}{分布函数}{distribution-function}
设 $X$ 是概率空间 $(\Omega,\mathscr F,P)$ 上的随机变量。函数
\[
F(t):=P(X\leq t),\qquad t\in\mathbb R,
\]
称为随机变量 $X$ 的\textbf{分布函数}。需要区分不同随机变量时，也记为 $F_X(t)$。
\end{mydefinition}
\origin{原笔记第2.1节，定义“分布函数”；源码第73行。}

\begin{mytheorem}{分布函数的基本性质}{distribution-function-properties}
设 $F$ 为随机变量 $X$ 的分布函数，则：
\begin{enumerate}
    \item $F$ 单调不减，即 $a\leq b$ 时有 $F(a)\leq F(b)$；
    \item $\displaystyle\lim_{t\to-\infty}F(t)=0$，$\displaystyle\lim_{t\to+\infty}F(t)=1$；
    \item $F$ 右连续，即对任意 $a\in\mathbb R$，
    \[
    F(a)=\lim_{t\to a^+}F(t);
    \]
    \item 对任意 $a\in\mathbb R$，记
    \[
    F(a^-):=\lim_{t\to a^-}F(t),
    \]
    则
    \[
    P(X=a)=F(a)-F(a^-).
    \]
\end{enumerate}
\end{mytheorem}
\origin{原笔记第2.1节，定理“分布函数的基本性质”；源码第85行。}

\subsection{例题：含跳跃的分布函数}\label{note:c2:e2}
\basis{原笔记第2.1节，例“含跳跃的分布函数”；源码\texttt{chapters/chapter2.tex}第169行。}
随机变量 $X$ 的分布函数为
\[
F(t)=
\begin{cases}
0,&t<0,\\[0.3em]
\dfrac t2,&0\leq t<1,\\[0.5em]
\dfrac23,&1\leq t<2,\\[0.5em]
\dfrac{11}{12},&2\leq t<3,\\[0.5em]
1,&t\geq3.
\end{cases}
\]
分别计算：
\[
P(X<3)=F(3^-)=\frac{11}{12},
\]
\[
P(X=1)=F(1)-F(1^-)=\frac23-\frac12=\frac16,
\]
\[
P\left(X>\frac12\right)=1-F\left(\frac12\right)
=1-\frac14=\frac34,
\]
\[
P(2<X\leq4)=F(4)-F(2)=1-\frac{11}{12}=\frac1{12},
\]
\[
P\left(\frac12\leq X<2\right)
=F(2^-)-F\left(\left(\frac12\right)^-\right)
=\frac23-\frac14=\frac5{12}.
\]

\subsection{例题：可列个分布函数的凸组合}\label{note:c2:e3}
\basis{原笔记第2.1节，例“可列个分布函数的凸组合”；源码\texttt{chapters/chapter2.tex}第204行。}
设 $\{F_m\}_{m\geq1}$ 是一列分布函数，$a_m\geq0$ 且
\[
\sum_{m=1}^{\infty}a_m=1.
\]
令
\[
F(t)=\sum_{m=1}^{\infty}a_mF_m(t).
\]
称 $F$ 为这些分布函数的\textbf{凸组合}。证明 $F$ 也是分布函数。

首先，
\[
0\leq\sum_{m=1}^{\infty}a_mF_m(t)
\leq\sum_{m=1}^{\infty}a_m=1,
\]
所以级数收敛，且 $F$ 定义良好。

（1）若 $a\leq b$，则每个 $F_m(a)\leq F_m(b)$，从而
\[
F(a)=\sum_{m=1}^{\infty}a_mF_m(a)
\leq\sum_{m=1}^{\infty}a_mF_m(b)=F(b).
\]

（2）固定 $a\in\mathbb R$ 和 $\varepsilon>0$。选择正整数 $M$，使
\[
\sum_{m=M+1}^{\infty}a_m<\frac{\varepsilon}{2}.
\]
由 $F_1,\ldots,F_M$ 右连续，存在 $\delta>0$，使 $a<t<a+\delta$ 时，对所有 $1\leq m\leq M$ 都有
\[
0\leq F_m(t)-F_m(a)<\frac{\varepsilon}{2}.
\]
因此
\[
\begin{aligned}
0\leq F(t)-F(a)
&=\sum_{m=1}^{M}a_m\bigl(F_m(t)-F_m(a)\bigr)\\
&\quad+\sum_{m=M+1}^{\infty}a_m\bigl(F_m(t)-F_m(a)\bigr)\\
&\leq\frac{\varepsilon}{2}\sum_{m=1}^{M}a_m
+\sum_{m=M+1}^{\infty}a_m\\
&<\varepsilon.
\end{aligned}
\]
故 $F$ 右连续。

（3）仍选取上述 $M$。由于有限个分布函数 $F_1,\ldots,F_M$ 在负无穷处均趋于 $0$，当 $t$ 足够小时，
\[
0\leq F(t)
\leq\frac{\varepsilon}{2}\sum_{m=1}^{M}a_m
+\sum_{m=M+1}^{\infty}a_m
<\varepsilon.
\]
故 $\lim_{t\to-\infty}F(t)=0$。同理，当 $t$ 足够大时，$1-F_m(t)<\varepsilon/2$ 对所有 $1\leq m\leq M$ 成立，于是
\[
0\leq1-F(t)
=\sum_{m=1}^{\infty}a_m\bigl(1-F_m(t)\bigr)
<\varepsilon.
\]
故 $\lim_{t\to+\infty}F(t)=1$。

综上，$F$ 满足分布函数的三个本质条件，因此 $F$ 是一个分布函数。

\clearpage\section{离散型随机变量}\label{k:discrete}

\begin{mydefinition}{离散型随机变量与分布列}{discrete-random-variable-pmf}
若随机变量 $X$ 的可能取值为有限个 $x_1,\ldots,x_s$ 或可列个 $x_1,x_2,\ldots$，则称 $X$ 为\textbf{离散型随机变量}。这些取值互不相同。令
\[
p(x_i):=P(X=x_i),
\]
称各取值及其对应概率构成的表为 $X$ 的\textbf{概率分布列}，简称\textbf{分布列}。函数 $p$ 也称为概率质量函数。

分布列可写成表格
\[
\begin{array}{c|ccccc}
X&x_1&x_2&\cdots&x_s&\cdots\\ \hline
P&p(x_1)&p(x_2)&\cdots&p(x_s)&\cdots
\end{array}
\]
或两行矩阵的形式
\[
X\sim
\begin{pmatrix}
x_1&x_2&\cdots&x_s&\cdots\\
p(x_1)&p(x_2)&\cdots&p(x_s)&\cdots
\end{pmatrix}.
\]
有限取值的情形去掉末尾的省略号。
\end{mydefinition}
\origin{原笔记第2.2节，定义“离散型随机变量与分布列”；源码第274行。}

\begin{proposition}{分布列的基本性质}{pmf-basic-properties}
离散型随机变量的分布列满足
\[
p(x_i)\geq0,\qquad \sum_i p(x_i)=1.
\]
其中求和遍历所有不同取值；有限情形为有限和，可列情形为无穷级数。
\end{proposition}
\origin{原笔记第2.2节，命题“分布列的基本性质”；源码第301行。}

\begin{proposition}{分布列与分布函数相互确定}{pmf-cdf-equivalence}
设 $X$ 为离散型随机变量，分布列为 $p(x_i)$，分布函数为 $F$。则
\[
F(t)=\sum_{i:x_i\leq t}p(x_i),\qquad t\in\mathbb R,
\]
并且
\[
p(x_i)=F(x_i)-F(x_i^-).
\]
\end{proposition}
\origin{原笔记第2.2节，命题“分布列与分布函数相互确定”；源码第326行。}

\subsection{例题：收集齐优惠券所需的张数}\label{note:c2:e4}
\basis{原笔记第2.2节，例“收集齐优惠券所需的张数”；源码\texttt{chapters/chapter2.tex}第341行。}
有 $r$ 种不同优惠券，每次收集一张，每一种被收集到的概率均为 $1/r$，各次收集相互独立。设 $T$ 为首次收集齐全部 $r$ 种优惠券时累计收集的张数，求 $P(T=n)$。

当 $r=1$ 时，$P(T=1)=1$。以下设 $r\geq2$；显然 $n<r$ 时，$P(T=n)=0$。

先求 $P(T>m)$。设 $A_j^{(m)}$ 表示前 $m$ 次收集中没有第 $j$ 种优惠券，则
\[
\{T>m\}=\bigcup_{j=1}^rA_j^{(m)}.
\]
对任意不同的 $j_1,\ldots,j_k$，由独立性，
\[
P\left(A_{j_1}^{(m)}\cap\cdots\cap A_{j_k}^{(m)}\right)
=\left(1-\frac kr\right)^m.
\]
当 $m\geq r-1\geq1$ 时，容斥公式给出
\[
\begin{aligned}
P(T>m)
&=\sum_{k=1}^r(-1)^{k+1}\binom rk
\left(1-\frac kr\right)^m\\
&=\sum_{k=1}^{r-1}(-1)^{k+1}\binom rk
\left(1-\frac kr\right)^m,
\end{aligned}
\]
其中 $k=r$ 对应的项为 $0$。于是对 $n\geq r$，
\[
\begin{aligned}
P(T=n)
&=P(T>n-1)-P(T>n)\\
&=\sum_{k=1}^{r-1}(-1)^{k+1}\binom rk
\left[
\left(1-\frac kr\right)^{n-1}
-\left(1-\frac kr\right)^n
\right].
\end{aligned}
\]

\begin{mydefinition}{二项分布}{binomial-distribution}
设 $n$ 为正整数，$0<p<1$。若随机变量 $X$ 取值于 $0,1,\ldots,n$，且
\[
P(X=k)=\binom nk p^k(1-p)^{n-k},
\qquad k=0,1,\ldots,n,
\]
则称 $X$ 服从参数为 $(n,p)$ 的\textbf{二项分布}，记为 $X\sim B(n,p)$。
\end{mydefinition}
\origin{原笔记第2.2节，定义“二项分布”；源码第381行。}

\begin{proposition}{二项分布的单调性与最可能取值}{binomial-modes}
设 $X\sim B(n,p)$，$0<p<1$，令 $a=(n+1)p$。随着 $k$ 从 $0$ 增加到 $n$，概率 $P(X=k)$ 先不减后不增，最大值的位置为：
\begin{enumerate}
    \item 若 $a$ 不是整数，则唯一的最可能取值为 $\lfloor a\rfloor$；
    \item 若 $a=m$ 是整数，则 $m-1$ 和 $m$ 是两个并列的最可能取值，且 $1\leq m\leq n$。
\end{enumerate}
其中 $\lfloor a\rfloor$ 表示不大于 $a$ 的最大整数。
\end{proposition}
\origin{原笔记第2.2节，命题“二项分布的单调性与最可能取值”；源码第407行。}

\subsection{例题：均匀硬币的中心项}\label{note:c2:e5}
\basis{原笔记第2.2节，例“均匀硬币的中心项”；源码\texttt{chapters/chapter2.tex}第436行。}
独立地抛一枚均匀硬币 $2n$ 次，设 $X_n$ 为正面出现的次数，则
\[
X_n\sim B\left(2n,\frac12\right).
\]
由于 $(2n+1)/2=n+1/2$ 不是整数，唯一的最可能正面次数为 $n$，且
\[
P(X_n=n)=\binom{2n}{n}\frac1{2^{2n}}
=\frac{(2n)!}{(n!)^2}\frac1{2^{2n}}.
\]
由斯特林（Stirling）公式
\[
n!\sim\sqrt{2\pi n}\left(\frac ne\right)^n,
\qquad n\to\infty,
\]
可得
\[
\begin{aligned}
P(X_n=n)
&\sim
\frac{\sqrt{4\pi n}(2n/e)^{2n}}
{2\pi n(n/e)^{2n}}\frac1{2^{2n}}\\
&=\frac1{\sqrt{\pi n}}\longrightarrow0.
\end{aligned}
\]
这里 $u_n\sim v_n$ 表示 $u_n/v_n\to1$。对所有 $0\leq k\leq2n$，
\[
0\leq P(X_n=k)\leq P(X_n=n).
\]
因此，对每个固定的非负整数 $k$，由夹逼定理，
\[
\lim_{n\to\infty}P(X_n=k)=0.
\]

\begin{mydefinition}{泊松分布}{poisson-distribution}
设 $\lambda>0$。若随机变量 $X$ 取值于非负整数，且
\[
P(X=k)=\frac{\lambda^k}{k!}e^{-\lambda},
\qquad k=0,1,2,\ldots,
\]
则称 $X$ 服从参数为 $\lambda$ 的\textbf{泊松分布}，记为 $X\sim\mathcal P(\lambda)$。
\end{mydefinition}
\origin{原笔记第2.2节，定义“泊松分布”；源码第473行。}

\subsection{例题：一页上至少有一处印刷错误}\label{note:c2:e6}
\basis{原笔记第2.2节，例“一页上至少有一处印刷错误”；源码\texttt{chapters/chapter2.tex}第492行。}
设一本书某一页上的印刷错误数 $X$ 服从参数为 $\lambda=1/2$ 的泊松分布。则
\[
P(X\geq1)=1-P(X=0)=1-e^{-1/2}\approx0.393469.
\]

\subsection{例题：泊松粒子数的独立捕捉}\label{note:c2:e7}
\basis{原笔记第2.2节，例“泊松粒子数的独立捕捉”；源码\texttt{chapters/chapter2.tex}第499行。}
设某仪器在一定时间内发射的粒子总数 $\xi\sim\mathcal P(\lambda)$。给定 $\xi=n$ 时，每个粒子以相同概率 $p$ 被捕捉记录，且各粒子的捕捉结果相互独立，其中 $0<p<1$。记捕捉到的粒子总数为 $\eta$，证明 $\eta\sim\mathcal P(\lambda p)$。

令 $q=1-p$。给定 $\xi=n$ 时，$\eta$ 的条件分布为 $B(n,p)$。对任意非负整数 $k$，由全概率公式，
\[
\begin{aligned}
P(\eta=k)
&=\sum_{n=k}^{\infty}P(\xi=n)P(\eta=k\mid\xi=n)\\
&=\sum_{n=k}^{\infty}\frac{\lambda^n}{n!}e^{-\lambda}
\binom nk p^kq^{n-k}\\
&=\frac{e^{-\lambda}(\lambda p)^k}{k!}
\sum_{n=k}^{\infty}\frac{(\lambda q)^{n-k}}{(n-k)!}\\
&=\frac{e^{-\lambda}(\lambda p)^k}{k!}
\sum_{j=0}^{\infty}\frac{(\lambda q)^j}{j!}\\
&=\frac{(\lambda p)^k}{k!}e^{-\lambda}e^{\lambda q}
=\frac{(\lambda p)^k}{k!}e^{-\lambda p}.
\end{aligned}
\]
其中换元 $j=n-k$，所以 $\eta\sim\mathcal P(\lambda p)$。

\begin{proposition}{泊松分布的单调性与最可能取值}{poisson-modes}
设 $X\sim\mathcal P(\lambda)$。随着 $k$ 从 $0$ 增大，$P(X=k)$ 先不减后不增，并趋于 $0$。
\begin{enumerate}
    \item 若 $\lambda$ 不是整数，则唯一的最可能取值为 $\lfloor\lambda\rfloor$；
    \item 若 $\lambda=m$ 为正整数，则 $m-1$ 和 $m$ 是两个并列的最可能取值。
\end{enumerate}
\end{proposition}
\origin{原笔记第2.2节，命题“泊松分布的单调性与最可能取值”；源码第520行。}

\begin{mytheorem}{泊松定理}{poisson-limit-theorem}
设 $X_n\sim B(n,p_n)$，$0<p_n<1$，且
\[
np_n\longrightarrow\lambda>0.
\]
则对每个固定的非负整数 $k$，
\[
\lim_{n\to\infty}P(X_n=k)=\frac{\lambda^k}{k!}e^{-\lambda}.
\]
\end{mytheorem}
\origin{原笔记第2.2节，定理“泊松定理”；源码第546行。}

\subsection{例题：死亡人数不超过五人的泊松近似}\label{note:c2:e8}
\basis{原笔记第2.2节，例“死亡人数不超过五人的泊松近似”；源码\texttt{chapters/chapter2.tex}第588行。}
某患病人群中，每人因该病死亡的概率为 $0.2\%$。在 $1000$ 人的结果相互独立且具有相同概率的模型下，设死亡人数为 $X$，则
\[
X\sim B(1000,0.002).
\]
精确概率为
\[
P(X\leq5)
=\sum_{k=0}^{5}\binom{1000}{k}0.002^k0.998^{1000-k}.
\]
令 $Z\sim\mathcal P(2)$，因为 $np=1000\times0.002=2$，所以
\[
P(X\leq5)\approx P(Z\leq5)
=e^{-2}\sum_{k=0}^{5}\frac{2^k}{k!}
\approx0.983436\approx0.9834.
\]
这是对原二项分布概率的近似，不能将 $X$ 的实际分布直接改写为泊松分布。

\begin{mydefinition}{几何分布}{geometric-distribution}
设 $0<p<1$。若随机变量 $X$ 取值于正整数，且
\[
P(X=k)=(1-p)^{k-1}p,\qquad k=1,2,\ldots,
\]
则称 $X$ 服从参数为 $p$ 的\textbf{几何分布}，记为 $X\sim G(p)$。
\end{mydefinition}
\origin{原笔记第2.2节，定义“几何分布”；源码第618行。}

\begin{mytheorem}{几何分布的无记忆性}{geometric-memorylessness}
设 $X$ 是取正整数值的随机变量，且对每个非负整数 $k$ 都有 $P(X>k)>0$。则 $X$ 服从某个参数 $0<p<1$ 的几何分布，当且仅当对任意正整数 $n$ 和非负整数 $k$，
\[
P(X=n+k\mid X>k)=P(X=n).
\]
\end{mytheorem}
\origin{原笔记第2.2节，定理“几何分布的无记忆性”；源码第642行。}

\begin{mydefinition}{帕斯卡分布}{pascal-distribution}
设 $n$ 为正整数，$0<p<1$。若随机变量 $X$ 取值于 $n,n+1,\ldots$，且
\[
P(X=k)=\binom{k-1}{n-1}p^n(1-p)^{k-n},
\qquad k=n,n+1,\ldots,
\]
则称 $X$ 服从参数为 $(n,p)$ 的\textbf{帕斯卡分布}，记为 $X\sim\mathrm{Pa}(n,p)$。
\end{mydefinition}
\origin{原笔记第2.2节，定义“帕斯卡分布”；源码第686行。}

\begin{mydefinition}{负二项分布}{negative-binomial-distribution}
若 $X\sim\mathrm{Pa}(n,p)$，令 $Y=X-n$，则 $Y$ 表示取得第 $n$ 次成功时已经发生的\textbf{失败次数}，取值于非负整数。其分布列为
\[
\begin{aligned}
P(Y=k)&=P(X=n+k)\\
&=\binom{n+k-1}{n-1}p^n(1-p)^k\\
&=\binom{n+k-1}{k}p^n(1-p)^k,\qquad k=0,1,\ldots.
\end{aligned}
\]
称 $Y$ 服从参数为 $(n,p)$ 的\textbf{负二项分布}，记为 $Y\sim\mathrm{NB}(n,p)$。
\end{mydefinition}
\origin{原笔记第2.2节，定义“负二项分布”；源码第715行。}

\subsection{例题：巴拿赫火柴盒问题}\label{note:c2:e9}
\basis{原笔记第2.2节，例“巴拿赫火柴盒问题”；源码\texttt{chapters/chapter2.tex}第739行。}
一个人左右口袋各放一盒火柴，每盒最初有 $N$ 根，$N$ 为正整数。每次需要火柴时，相互独立地以概率 $1/2$ 选择一盒；若盒内有火柴，就取出一根。求他首次选到一个空盒、发现该盒已用光时，另一盒恰好剩下 $k$ 根的概率，其中 $0\leq k\leq N$。

设 $A_k$ 为首次发现右盒为空时，左盒恰好剩下 $k$ 根的事件。由于左右对称，两种首次发现空盒的事件互斥，所求概率为 $2P(A_k)$。

$A_k$ 发生时，这是第 $N+1$ 次选择右盒，而此前共选择左盒 $N-k$ 次。将选择右盒视为成功，则
\[
T\sim\mathrm{Pa}\left(N+1,\frac12\right),
\]
其中 $T$ 表示第 $N+1$ 次选择右盒的时刻。$A_k$ 等价于 $T=2N-k+1$；此前选择左盒的次数不超过 $N$，因而还没有发现左盒为空。于是
\[
\begin{aligned}
P(A_k)
&=\binom{2N-k}{N}
\left(\frac12\right)^{N+1}
\left(\frac12\right)^{N-k}\\
&=\binom{2N-k}{N}\frac1{2^{2N-k+1}}.
\end{aligned}
\]
因此所求概率为
\[
\boxed{\,P(\text{另一盒剩下 }k\text{ 根})
=\binom{2N-k}{N}\frac1{2^{2N-k}},\quad 0\leq k\leq N.\,}
\]

\textbf{注：}\quad 停止时刻是“首次选到空盒并发现已经用光”，不是“刚取走某盒的最后一根火柴”；区分这两个时刻才能得到正确的分布。

\clearpage\section{连续型随机变量}\label{k:continuous}

\begin{mydefinition}{连续型随机变量与概率密度}{continuous-random-variable-density}
设 $X$ 为随机变量。若存在非负可测函数 $f:\mathbb R\to[0,\infty)$，使得对任意 $B\in\mathscr B(\mathbb R)$，
\[
P(X\in B)=\int_B f(x)\,dx,
\]
则称 $X$ 为\textbf{连续型随机变量}，称 $f$ 为 $X$ 的\textbf{概率密度函数}，简称\textbf{密度函数}或\textbf{密度}。
\end{mydefinition}
\origin{原笔记第2.3节，定义“连续型随机变量与概率密度”；源码第794行。}

\begin{proposition}{概率密度的基本性质}{density-basic-properties}
若 $f$ 是连续型随机变量 $X$ 的密度，则
\[
f(x)\geq0,\qquad
\int_{\mathbb R}f(x)\,dx
=\int_{-\infty}^{+\infty}f(x)\,dx=1.
\]
\end{proposition}
\origin{原笔记第2.3节，命题“概率密度的基本性质”；源码第804行。}

\subsection{例题：分段密度的参数、区间概率与分布函数}\label{note:c2:e10}
\basis{原笔记第2.3节，例“分段密度的参数、区间概率与分布函数”；源码\texttt{chapters/chapter2.tex}第828行。}
设 $X$ 的概率密度为
\[
f(x)=
\begin{cases}
\dfrac18,&0<x<2,\\[0.4em]
kx,&2\leq x<4,\\
0,&\text{其他}.
\end{cases}
\]
求 $k$、$P(1<X\leq3)$ 以及分布函数 $F$。

（1）由密度的规范性，
\[
1=\int_0^2\frac18\,dx+\int_2^4kx\,dx
=\frac14+6k,
\]
故 $k=1/8$，也满足非负性要求。

（2）区间跨越密度的分段点 $2$，因此
\[
P(1<X\leq3)
=\int_1^2\frac18\,dx+\int_2^3\frac{x}{8}\,dx
=\frac18+\frac5{16}=\frac7{16}.
\]

（3）由 $F(t)=\int_{-\infty}^t f(x)\,dx$ 分段计算。特别地，$2\leq t<4$ 时，
\[
F(t)=\int_0^2\frac18\,dx+\int_2^t\frac{x}{8}\,dx
=\frac14+\frac{t^2-4}{16}=\frac{t^2}{16}.
\]
综上，
\[
F(t)=
\begin{cases}
0,&t<0,\\[0.3em]
\dfrac t8,&0\leq t<2,\\[0.5em]
\dfrac{t^2}{16},&2\leq t<4,\\[0.5em]
1,&t\geq4.
\end{cases}
\]
在 $0,2,4$ 处左右极限与函数值均相同，故 $F$ 连续。

\textbf{注：}\quad “几乎处处”是指除去一个勒贝格测度为 $0$ 的集合后成立；这里的测度是区间长度概念的推广。改变密度在单个点或零测集上的非负取值，不改变其所描述的分布。求导等式须按上述条件和几乎处处的意义理解。

\begin{proposition}{连续型分布的基本性质}{continuous-distribution-properties}
设 $X$ 为连续型随机变量，密度与分布函数分别为 $f$、$F$。则：
\begin{enumerate}
    \item 对任意 $t\in\mathbb R$，
    \[
    F(t)=\int_{-\infty}^t f(x)\,dx;
    \]
    \item $F$ 连续，且几乎处处可微，并且
    \[
    F'(t)=f(t)\qquad\text{几乎处处成立};
    \]
    在 $f$ 连续的点 $t$，由积分上限函数求导可得 $F'(t)=f(t)$；
    \item 若 Borel 集 $B$ 的勒贝格测度 $m(B)=0$，则
    \[
    P(X\in B)=0.
    \]
    特别地，对任意实数 $c$，$P(X=c)=0$；
    \item 对给定的随机变量 $X$，分布函数唯一，密度则只在相差一个零测集的意义下唯一。即若 $f_1,f_2$ 都是其密度，则
    \[
    m\bigl(\{x:f_1(x)\ne f_2(x)\}\bigr)=0.
    \]
\end{enumerate}
\end{proposition}
\origin{原笔记第2.3节，命题“连续型分布的基本性质”；源码第874行。}

\subsection{例题：五个电子管中恰有两个需要更换}\label{note:c2:e11}
\basis{原笔记第2.3节，例“五个电子管中恰有两个需要更换”；源码\texttt{chapters/chapter2.tex}第934行。}
某类电子管的寿命 $X$（单位：小时）具有密度
\[
f(x)=
\begin{cases}
\dfrac{100}{x^2},&x>100,\\[0.4em]
0,&x\leq100.
\end{cases}
\]
现有 $5$ 个具有相同寿命分布且相互独立的电子管，求 $150$ 小时内恰有 $2$ 个原装电子管因寿命耗尽需要更换的概率。

单个电子管在 $150$ 小时内需要更换的概率为
\[
p=P(X\leq150)
=\int_{100}^{150}\frac{100}{x^2}\,dx
=\left[-\frac{100}{x}\right]_{100}^{150}
=\frac13.
\]
记需要更换的原装电子管个数为 $Y$。由寿命的独立性，
\[
Y\sim B\left(5,\frac13\right),
\]
所以
\[
P(Y=2)=\binom52\left(\frac13\right)^2
\left(\frac23\right)^3=\frac{80}{243}.
\]

\subsection{例题：由分布函数求密度}\label{note:c2:e12}
\basis{原笔记第2.3节，例“由分布函数求密度”；源码\texttt{chapters/chapter2.tex}第963行。}
设连续型随机变量 $X$ 的分布函数为
\[
F(t)=
\begin{cases}
0,&t<0,\\
t^2,&0\leq t<1,\\
1,&t\geq1.
\end{cases}
\]
在各个开区间上求导，可选择一个密度版本为
\[
f(x)=
\begin{cases}
2x,&0<x<1,\\
0,&\text{其他}.
\end{cases}
\]
因为 $F$ 在 $1$ 处不可微，$f(1)=0$ 是选定的密度值，不是该点的导数；零测集上的取值可以另作非负指定。并且
\[
\int_0^1 2x\,dx=1,
\]
将此密度从 $-\infty$ 积分到 $t$，恰好还原上面的 $F(t)$。

\textbf{注：}\quad 可简写为“$f(x)=2x,\ 0<x<1$，其余处为 $0$”；不能把区间外的零值条件省略成对所有实数都有 $f(x)=2x$。

\begin{mydefinition}{均匀分布}{uniform-distribution}
设 $a<b$。若随机变量 $X$ 的密度可以取为
\[
f(x)=
\begin{cases}
\dfrac1{b-a},&a<x<b,\\[0.4em]
0,&\text{其他},
\end{cases}
\]
则称 $X$ 服从区间 $(a,b)$ 上的\textbf{均匀分布}，记为 $X\sim U(a,b)$。
\end{mydefinition}
\origin{原笔记第2.3节，定义“均匀分布”；源码第1005行。}

\subsection{例题：等公交车不超过五分钟}\label{note:c2:e13}
\basis{原笔记第2.3节，例“等公交车不超过五分钟”；源码\texttt{chapters/chapter2.tex}第1031行。}
公交车从 $7$ 点起每隔 $15$ 分钟到站一次。乘客到站时间服从 $7$ 点至 $7$ 点 $30$ 分之间的均匀分布，求等车时间不超过 $5$ 分钟的概率。

设 $X$ 为乘客到站时距 $7$ 点的分钟数，则 $X\sim U(0,30)$。在本时间段内公交车于 $0,15,30$ 分钟到站，所求事件为
\[
A=\{X=0\}\cup\{10\leq X\leq15\}\cup\{25\leq X\leq30\}.
\]
上述事件互斥，且 $P(X=0)=0$，因此
\[
\begin{aligned}
P(A)
&=P(X=0)+P(10\leq X\leq15)+P(25\leq X\leq30)\\
&=0+\int_{10}^{15}\frac1{30}\,dx+\int_{25}^{30}\frac1{30}\,dx\\
&=\frac5{30}+\frac5{30}=\frac13.
\end{aligned}
\]

\begin{mydefinition}{指数分布}{exponential-distribution}
设 $\lambda>0$。若随机变量 $X$ 的密度可以取为
\[
f(x)=
\begin{cases}
\lambda e^{-\lambda x},&x\geq0,\\
0,&x<0,
\end{cases}
\]
则称 $X$ 服从参数为 $\lambda$ 的\textbf{指数分布}，记为 $X\sim\operatorname{Exp}(\lambda)$。这里 $\lambda$ 使用率参数的约定。
\end{mydefinition}
\origin{原笔记第2.3节，定义“指数分布”；源码第1051行。}

\begin{proposition}{指数分布的分布函数与尾概率}{exponential-cdf}
对连续型随机变量 $X$，$X\sim\operatorname{Exp}(\lambda)$ 当且仅当
\[
F(t)=
\begin{cases}
1-e^{-\lambda t},&t\geq0,\\
0,&t<0.
\end{cases}
\]
等价地，其尾概率为
\[
P(X>t)=
\begin{cases}
e^{-\lambda t},&t\geq0,\\
1,&t<0.
\end{cases}
\]
\end{proposition}
\origin{原笔记第2.3节，命题“指数分布的分布函数与尾概率”；源码第1073行。}

\begin{mytheorem}{指数分布的无记忆性}{exponential-memorylessness}
设 $X$ 为取值非负的连续型随机变量，且对每个 $s>0$，$P(X>s)>0$。则 $X$ 服从某个参数 $\lambda>0$ 的指数分布，当且仅当对任意 $s,t>0$，
\[
P(X>s+t\mid X>s)=P(X>t).
\]
这一性质也称为\textbf{无效性}。
\end{mytheorem}
\origin{原笔记第2.3节，定理“指数分布的无记忆性”；源码第1108行。}

\begin{mylemma}{连续乘法方程}{continuous-multiplicative-equation}
若连续函数 $G:(0,\infty)\to(0,\infty)$ 满足
\[
G(s+t)=G(s)G(t),\qquad s,t>0,
\]
则存在 $b\in\mathbb R$，使 $G(t)=e^{bt}$ 对所有 $t>0$ 成立。
\end{mylemma}
\origin{原笔记第2.3节，引理“连续乘法方程”；源码第1151行。}

\subsection{例题：邮局中最后办完事情的人}\label{note:c2:e14}
\basis{原笔记第2.3节，例“邮局中最后办完事情的人”；源码\texttt{chapters/chapter2.tex}第1200行。}
邮局有两位职员。$A$ 到达时，$B$、$C$ 正分别接受服务；其中一人办完后，$A$ 立即开始接受服务。假定各人的服务时间相互独立，均服从同一 $\operatorname{Exp}(\lambda)$ 分布，求 $A$ 在三人中最后办完的概率。

由无记忆性，$A$ 开始服务时，另一位尚未办完者的剩余服务时间 $R$ 仍服从 $\operatorname{Exp}(\lambda)$。$A$ 的服务时间 $T$ 也服从同一分布，且与 $R$ 独立。此时已有一人办完，因此
\[
\{A\text{ 最后办完}\}=\{T>R\}.
\]
独立同分布给出对称性，
\[
P(T>R)=P(R>T).
\]
两者独立且为连续型，$P(T=R)=0$，所以
\[
P(T>R)+P(R>T)=1,
\qquad
P(A\text{ 最后办完})=\frac12.
\]

\textbf{注：}\quad 该模型要求服务时间相互独立；仅有相同的指数边缘分布，并不足以推出上述比较结论。

\subsection{例题：第一位顾客的到达时刻}\label{note:c2:e15}
\basis{原笔记第2.3节，例“第一位顾客的到达时刻”；源码\texttt{chapters/chapter2.tex}第1221行。}
设 $N(t)$ 表示商场从开门时刻 $0$ 到时刻 $t$ 的顾客到达人数，且对每个 $t>0$，
\[
N(t)\sim\mathcal P(\lambda t),\qquad \lambda>0.
\]
令 $S$ 为第一位顾客的到达时刻。对 $t>0$，
\[
\{S>t\}=\{N(t)=0\},
\]
所以
\[
\begin{aligned}
P(S>t)
&=P(N(t)=0)\\
&=e^{-\lambda t}.
\end{aligned}
\]
结合到达时刻非负可知，$S\sim\operatorname{Exp}(\lambda)$。

\begin{mydefinition}{Gamma 函数与 Gamma 分布}{gamma-function-and-distribution}
对 $\alpha>0$，定义
\[
\Gamma(\alpha)=\int_0^\infty x^{\alpha-1}e^{-x}\,dx,
\]
称为\textbf{Gamma 函数}或 $\Gamma$ 函数。

设 $\alpha>0$、$\lambda>0$。若随机变量 $X$ 的密度可以取为
\[
f(x)=
\begin{cases}
\dfrac{\lambda^\alpha}{\Gamma(\alpha)}x^{\alpha-1}e^{-\lambda x},
&x>0,\\[0.5em]
0,&x\leq0,
\end{cases}
\]
则称 $X$ 服从参数为 $(\alpha,\lambda)$ 的\textbf{Gamma 分布}，记为 $X\sim\Gamma(\alpha,\lambda)$。其中 $\alpha$ 为形状参数，$\lambda$ 为率参数；正半轴上的密度也可写成
\[
f(x)=\frac{(\lambda x)^{\alpha-1}}{\Gamma(\alpha)}
\lambda e^{-\lambda x},\qquad x>0.
\]
\end{mydefinition}
\origin{原笔记第2.3节，定义“Gamma 函数与 Gamma 分布”；源码第1243行。}

\begin{proposition}{Gamma 函数的基本性质}{gamma-function-properties}
Gamma 函数满足：
\begin{enumerate}
    \item 对 $\alpha>0$，$\Gamma(\alpha+1)=\alpha\Gamma(\alpha)$；
    \item $\Gamma(1)=1$，$\Gamma(1/2)=\sqrt\pi$；
    \item 对正整数 $n$，$\Gamma(n)=(n-1)!$。
\end{enumerate}
\end{proposition}
\origin{原笔记第2.3节，命题“Gamma 函数的基本性质”；源码第1280行。}

\begin{proposition}{Gamma 密度的归一化}{gamma-density-normalization}
上述 Gamma 分布定义中的 $f$ 确实是概率密度函数。
\end{proposition}
\origin{原笔记第2.3节，命题“Gamma 密度的归一化”；源码第1363行。}

\subsection{例题：\texorpdfstring{第 $n$ 位顾客的到达时刻}{第 n 位顾客的到达时刻}}\label{note:c2:e16}
\basis{原笔记第2.3节，例“第 $n$ 位顾客的到达时刻”；源码\texttt{chapters/chapter2.tex}第1381行。}
仍设商场在 $[0,t]$ 内的到达人数满足
\[
N(t)\sim\mathcal P(\lambda t),\qquad t>0.
\]
令 $S_n$ 为第 $n$ 位顾客的到达时刻，$n\geq1$。证明 $S_n\sim\Gamma(n,\lambda)$。

先求分布函数，再求导。对 $t>0$，
\[
\{S_n\leq t\}=\{N(t)\geq n\}.
\]
所以
\[
\begin{aligned}
F_{S_n}(t)
&=P(N(t)\geq n)\\
&=1-P(N(t)\leq n-1)\\
&=1-e^{-\lambda t}\sum_{k=0}^{n-1}\frac{(\lambda t)^k}{k!}.
\end{aligned}
\]
$t\leq0$ 时取 $F_{S_n}(t)=0$，与右侧在 $0$ 处的极限一致。

对 $t>0$ 求导，将乘积的两部分分别列出：
\[
\begin{aligned}
f_{S_n}(t)
&=e^{-\lambda t}\sum_{k=0}^{n-1}
\frac{\lambda^{k+1}t^k}{k!}
-e^{-\lambda t}\sum_{k=1}^{n-1}
\frac{\lambda^kt^{k-1}}{(k-1)!}\\
&=e^{-\lambda t}\left[
\sum_{k=0}^{n-1}\frac{\lambda^{k+1}t^k}{k!}
-\sum_{j=0}^{n-2}\frac{\lambda^{j+1}t^j}{j!}
\right]\\
&=\frac{\lambda^nt^{n-1}}{(n-1)!}e^{-\lambda t}\\
&=\frac{(\lambda t)^{n-1}}{\Gamma(n)}
\lambda e^{-\lambda t}.
\end{aligned}
\]
前两式中两个求和的共有项逐项抵消；$n=1$ 时第二个求和为空和，按 $0$ 处理。取 $t\leq0$ 时密度为 $0$，便得到 Gamma 密度，故
\[
S_n\sim\Gamma(n,\lambda).
\]
特别地，$S_1\sim\Gamma(1,\lambda)=\operatorname{Exp}(\lambda)$，与第一位顾客的结果一致。

\begin{mydefinition}{贝塔函数与贝塔分布}{beta-function-and-distribution}
设 $a,b>0$，定义\textbf{贝塔函数}
\[
B(a,b)=\int_0^1x^{a-1}(1-x)^{b-1}\,dx.
\]
若随机变量 $X$ 的密度可以取为
\[
f(x)=
\begin{cases}
\dfrac{x^{a-1}(1-x)^{b-1}}{B(a,b)},&0<x<1,\\[0.5em]
0,&\text{其他},
\end{cases}
\]
则称 $X$ 服从参数为 $(a,b)$ 的\textbf{贝塔分布}。
\end{mydefinition}
\origin{原笔记第2.3节，定义“贝塔函数与贝塔分布”；源码第1429行。}

\begin{proposition}{贝塔函数的基本性质}{beta-function-properties}
对 $a,b>0$，
\[
B(a,b)=B(b,a),
\qquad
B(a,b)=\frac{\Gamma(a)\Gamma(b)}{\Gamma(a+b)}.
\]
\end{proposition}
\origin{原笔记第2.3节，命题“贝塔函数的基本性质”；源码第1452行。}

\begin{mydefinition}{正态分布}{normal-distribution}
设 $\mu\in\mathbb R$、$\sigma>0$。若随机变量 $X$ 的密度为
\[
f(x)=\frac1{\sqrt{2\pi}\sigma}
\exp\left(-\frac{(x-\mu)^2}{2\sigma^2}\right),
\qquad x\in\mathbb R,
\]
则称 $X$ 服从参数为 $(\mu,\sigma^2)$ 的\textbf{正态分布}或\textbf{高斯分布}，记为 $X\sim N(\mu,\sigma^2)$。特别地，$N(0,1)$ 称为\textbf{标准正态分布}。
\end{mydefinition}
\origin{原笔记第2.3节，定义“正态分布”；源码第1504行。}

\begin{proposition}{正态密度的归一化}{normal-density-normalization}
上述 $f$ 非负且在实数轴上的积分为 $1$，因而确实是概率密度。
\end{proposition}
\origin{原笔记第2.3节，命题“正态密度的归一化”；源码第1514行。}

\begin{proposition}{正态密度的图像性质}{normal-density-shape}
正态密度关于 $x=\mu$ 对称，在 $\mu$ 处唯一达到最大值 $1/(\sqrt{2\pi}\sigma)$；以 $x$ 轴为水平渐近线，拐点为 $x=\mu\pm\sigma$，图像呈钟形。

固定 $\sigma$ 改变 $\mu$ 时，图像水平平移；固定 $\mu$ 增大 $\sigma$ 时，图像展宽、峰值降低，减小 $\sigma$ 时相反。
\end{proposition}
\origin{原笔记第2.3节，命题“正态密度的图像性质”；源码第1535行。}

标准正态分布的密度和分布函数分别记为
\[
\varphi(x)=\frac1{\sqrt{2\pi}}e^{-x^2/2},
\qquad
\Phi(t)=\int_{-\infty}^t\varphi(x)\,dx.
\]

\begin{proposition}{标准正态分布的对称公式}{standard-normal-symmetry}
对任意 $x\in\mathbb R$，
\[
\Phi(x)+\Phi(-x)=1.
\]
\end{proposition}
\origin{原笔记第2.3节，命题“标准正态分布的对称公式”；源码第1577行。}

\subsection{例题：对两侧尾部作符号翻转}\label{note:c2:e17}
\basis{原笔记第2.3节，例“对两侧尾部作符号翻转”；源码\texttt{chapters/chapter2.tex}第1600行。}
设 $X\sim N(0,1)$，$a>0$，定义
\[
Y=
\begin{cases}
X,&|X|\leq a,\\
-X,&|X|>a.
\end{cases}
\]
证明 $Y\sim N(0,1)$。

先写出分布函数：
\[
\begin{aligned}
F_Y(y)
&=P(X\leq y,\ |X|\leq a)
+P(-X\leq y,\ |X|>a)\\
&=P(X\leq y,\ |X|\leq a)
+P(X\geq-y,\ |X|>a).
\end{aligned}
\]
分三种情况计算。

（1）若 $y>a$，三个相应区间互不相交，
\[
\begin{aligned}
F_Y(y)
&=P(|X|\leq a)+P(-y\leq X<-a)+P(X>a)\\
&=P(X\geq-y)=\Phi(y).
\end{aligned}
\]

（2）若 $-a\leq y\leq a$，利用两侧尾概率相等，
\[
\begin{aligned}
F_Y(y)
&=P(-a\leq X\leq y)+P(X>a)\\
&=P(-a\leq X\leq y)+P(X<-a)\\
&=P(X\leq y)=\Phi(y).
\end{aligned}
\]

（3）若 $y<-a$，中间区间不贡献概率，
\[
F_Y(y)=0+P(X\geq-y)=\Phi(y).
\]
综上 $F_Y=\Phi$，故 $Y\sim N(0,1)$。

\begin{mytheorem}{正态分布的线性性质}{normal-affine-transformation}
若 $X\sim N(\mu,\sigma^2)$，$a,b\in\mathbb R$ 且 $a\ne0$，则
\[
aX+b\sim N(a\mu+b,a^2\sigma^2).
\]
\end{mytheorem}
\origin{原笔记第2.3节，定理“正态分布的线性性质”；源码第1651行。}

\begin{mycorollary}{正态分布的标准化}{normal-standardization}
若 $X\sim N(\mu,\sigma^2)$，则
\[
\frac{X-\mu}{\sigma}\sim N(0,1).
\]
\end{mycorollary}
\origin{原笔记第2.3节，推论“正态分布的标准化”；源码第1700行。}

\subsection{例题：三倍尺度区间的概率}\label{note:c2:e18}
\basis{原笔记第2.3节，例“三倍尺度区间的概率”；源码\texttt{chapters/chapter2.tex}第1727行。}
设 $X\sim N(\mu,\sigma^2)$，则
\[
\begin{aligned}
P(|X-\mu|<3\sigma)
&=P\left(-3<\frac{X-\mu}{\sigma}<3\right)\\
&=\Phi(3)-\Phi(-3)\\
&=2\Phi(3)-1\approx0.9973002.
\end{aligned}
\]
若先取 $\Phi(3)\approx0.9987$，再计算得到 $0.9974$，这是粗查表近似；较精细的概率约为 $99.73\%$。

\textbf{注：}\quad 这称为 $3\sigma$ 原则，可用于理解质量控制中的异常提示。但区间外仍有约 $0.27\%$ 的概率，并非概率为 $0$；一次观测落在区间外，不能单凭这一点就断言随机变量不服从该正态分布。

\subsection{例题：由两端比例估计录取分数线}\label{note:c2:e19}
\basis{原笔记第2.3节，例“由两端比例估计录取分数线”；源码\texttt{chapters/chapter2.tex}第1742行。}
共有 $10000$ 人报考，成绩大致服从正态分布；$90$ 分以上有 $359$ 人，$60$ 分以下有 $1151$ 人，现从高到低录取 $2500$ 人，估计最低录取分数。

设成绩模型为 $X\sim N(\mu,\sigma^2)$，将人数比例近似为概率。首先，
\[
\begin{aligned}
P(X\geq90)&\approx\frac{359}{10000}=0.0359,\\
\Phi\left(\frac{90-\mu}{\sigma}\right)&\approx0.9641.
\end{aligned}
\]
查表得到
\[
\frac{90-\mu}{\sigma}\approx1.8.
\]
另一方面，
\[
\begin{aligned}
\Phi\left(\frac{60-\mu}{\sigma}\right)
&\approx0.1151,\\
\Phi\left(\frac{\mu-60}{\sigma}\right)
&\approx1-0.1151=0.8849.
\end{aligned}
\]
所以 $(\mu-60)/\sigma\approx1.2$。联立两式，
\[
\begin{cases}
90-\mu\approx1.8\sigma,\\
\mu-60\approx1.2\sigma,
\end{cases}
\qquad
\sigma\approx10,\quad \mu\approx72.
\]

设最低录取分为 $c$。取前 $25\%$，意味着
\[
\begin{aligned}
P(X\geq c)&\approx0.25,\\
\Phi\left(\frac{c-72}{10}\right)&\approx0.75.
\end{aligned}
\]
按粗查表值 $(c-72)/10\approx0.68$，得到 $c\approx78.8$ 分。若用较精细的 $0.75$ 分位数 $0.67449$，在同一 $\mu\approx72,\sigma\approx10$ 的近似下，
\[
c\approx72+10\times0.67449\approx78.74\text{ 分}.
\]
$78.8$ 分是粗查表结果，不是精确分位数。

\clearpage\section{随机变量函数的分布}\label{k:transform}

\begin{mydefinition}{博雷尔函数}{borel-function}
若函数 $g:\mathbb R\to\mathbb R$ 满足
\[
g^{-1}(B)\in\mathcal B(\mathbb R),
\qquad \forall B\in\mathcal B(\mathbb R),
\]
则称 $g$ 为\textbf{博雷尔函数}，也称 $(\mathbb R,\mathcal B(\mathbb R))$ 上的可测函数。
\end{mydefinition}
\origin{原笔记第2.4节，定义“博雷尔函数”；源码第1803行。}

\subsection{例题：离散型随机变量的多项式变换}\label{note:c2:e20}
\basis{原笔记第2.4节，例“离散型随机变量的多项式变换”；源码\texttt{chapters/chapter2.tex}第1818行。}
已知 $X$ 的分布列为
\[
\begin{array}{c|ccccc}
X&-2&-1&0&1&2\\ \hline
P&0.2&0.1&0.1&0.3&0.3
\end{array},
\]
求 $Y=X^2+X$ 的分布列。

逐项计算并保留原来的概率：
\[
\begin{array}{c|ccccc}
X&-2&-1&0&1&2\\ \hline
Y=X^2+X&2&0&0&2&6\\
P&0.2&0.1&0.1&0.3&0.3
\end{array}.
\]
$Y=0$ 对应 $X=-1,0$，$Y=2$ 对应 $X=-2,1$，$Y=6$ 对应 $X=2$。合并后得到
\[
\begin{array}{c|ccc}
Y&0&2&6\\ \hline
P&0.2&0.5&0.3
\end{array}.
\]
三个概率之和为 $1$。

\subsection{例题：一般连续型随机变量的线性变换}\label{note:c2:e21}
\basis{原笔记第2.4节，例“一般连续型随机变量的线性变换”；源码\texttt{chapters/chapter2.tex}第1848行。}
设 $X$ 的密度为 $f_X$，求 $Y=aX+b$（$a\ne0$）的密度。

先按 $a$ 的符号求分布函数。$a>0$ 时，
\[
\begin{aligned}
F_Y(t)
&=P\left(X\leq\frac{t-b}{a}\right)\\
&=F_X\left(\frac{t-b}{a}\right).
\end{aligned}
\]
$a<0$ 时，利用 $X$ 无单点质量，
\[
\begin{aligned}
F_Y(t)
&=P\left(X\geq\frac{t-b}{a}\right)\\
&=1-F_X\left(\frac{t-b}{a}\right).
\end{aligned}
\]
分别对上式求导，合并得到一个密度版本
\[
\boxed{\quad f_Y(t)=\frac1{|a|}f_X\left(\frac{t-b}{a}\right).\quad}
\]

\begin{mydefinition}{分布函数的广义逆}{generalized-inverse-cdf}
设 $G$ 单调不减、右连续，且
\[
\lim_{x\to-\infty}G(x)=0,\qquad
\lim_{x\to+\infty}G(x)=1.
\]
对 $p\in(0,1)$，定义其\textbf{广义逆}
\[
G^{-1}(p)=\inf\{x\in\mathbb R:G(x)\geq p\}.
\]
这里不要求 $G$ 严格单调，也不要求它有通常意义的反函数。
\end{mydefinition}
\origin{原笔记第2.4节，定义“分布函数的广义逆”；源码第1887行。}

\begin{mylemma}{广义逆与分布函数的不等式}{generalized-inverse-equivalence}
对任意 $p\in(0,1)$ 和 $t\in\mathbb R$，
\[
G^{-1}(p)\leq t\quad\Longleftrightarrow\quad p\leq G(t).
\]
\end{mylemma}
\origin{原笔记第2.4节，引理“广义逆与分布函数的不等式”；源码第1906行。}

\begin{mytheorem}{由均匀分布构造给定分布}{inverse-transform-construction}
设 $X\sim U(0,1)$，$G$ 及其广义逆满足前述定义。令
\[
Y=G^{-1}(X),
\]
则 $Y$ 的分布函数就是 $G$。
\end{mytheorem}
\origin{原笔记第2.4节，定理“由均匀分布构造给定分布”；源码第1935行。}

\begin{mytheorem}{分支反函数的密度变换公式}{inverse-branches-density}
设 $X$ 有密度 $f_X$，$g$ 为博雷尔函数，$Y=g(X)$。设存在有限个或可数个互不相交的开区间 $I_i$，满足
\[
P\left(X\in\bigcup_i I_i\right)=1.
\]
在每个 $I_i$ 上，$g$ 连续、严格单调，像 $J_i=g(I_i)$ 为开区间；其反函数 $h_i:J_i\to I_i$ 连续可微。这里 $h_i$ 称为第 $i$ 个\textbf{反函数分支}，即
\[
g(h_i(y))=y,\qquad y\in J_i.
\]
则 $Y$ 有一个密度版本
\[
f_Y(y)=\sum_{\{i:y\in J_i\}}f_X(h_i(y))\,|h_i'(y)|.
\]
没有分支贡献时，该和取 $0$。非负级数在几乎处处有限；在可能出现的零测异常点上，可另赋有限值而不改变分布。
\end{mytheorem}
\origin{原笔记第2.4节，定理“分支反函数的密度变换公式”；源码第1977行。}

\subsection{例题：用反函数公式求线性变换的密度}\label{note:c2:e22}
\basis{原笔记第2.4节，例“用反函数公式求线性变换的密度”；源码\texttt{chapters/chapter2.tex}第1993行。}
设 $Y=aX+b$，$a\ne0$。变换在整个实轴上只有一个反函数分支
\[
h(y)=\frac{y-b}{a},\qquad |h'(y)|=\frac1{|a|}.
\]
所以
\[
f_Y(y)=\frac1{|a|}f_X\left(\frac{y-b}{a}\right).
\]
这与分布函数法的结果一致。

\subsection{例题：平方变换与不同分支的取值范围}\label{note:c2:e23}
\basis{原笔记第2.4节，例“平方变换与不同分支的取值范围”；源码\texttt{chapters/chapter2.tex}第2005行。}
设随机变量 $X$ 的密度为
\[
f_X(x)=
\begin{cases}
\dfrac14|x|,&-2\leq x<0,\\
x,&0\leq x\leq1,\\
0,&\text{其他}.
\end{cases}
\]
求 $Y=X^2$ 的密度。

负半轴的概率为 $\int_{-2}^0|x|/4\,dx=1/2$，正半轴的概率为 $\int_0^1x\,dx=1/2$，所以给定密度归一化。由 $-2\leq X\leq1$ 以概率 $1$ 成立，$Y$ 可限制在 $[0,4]$。

对 $y>0$，两个反函数分支为 $h_1(y)=\sqrt y$ 和 $h_2(y)=-\sqrt y$，它们的导数绝对值都为 $1/(2\sqrt y)$，因此
\[
f_Y(y)=\frac{f_X(\sqrt y)+f_X(-\sqrt y)}{2\sqrt y}.
\]
当 $0<y<1$ 时，正、负两个原像均落在原密度的非零区间内，故
\[
f_Y(y)=\frac{\sqrt y+\sqrt y/4}{2\sqrt y}=\frac58.
\]
当 $1<y<4$ 时，正原像超过 $1$，只有负分支贡献，故
\[
f_Y(y)=\frac{0+\sqrt y/4}{2\sqrt y}=\frac18.
\]
当 $y<0$ 或 $y>4$ 时，概率密度可取 $0$。在 $y=0$ 处不能代入含 $1/\sqrt y$ 的公式；但
\[
P(Y=0)=P(X=0)=0,
\]
可另指定该点的密度值。取密度版本为
\[
f_Y(y)=
\begin{cases}
\dfrac58,&0\leq y\leq1,\\
\dfrac18,&1<y\leq4,\\
0,&\text{其他}.
\end{cases}
\]
检查全积分：
\[
\int_{\mathbb R}f_Y(y)\,dy=\frac58+\frac18(4-1)=1.
\]

\subsection{例题：均匀随机变量的余弦变换}\label{note:c2:e24}
\basis{原笔记第2.4节，例“均匀随机变量的余弦变换”；源码\texttt{chapters/chapter2.tex}第2050行。}
设 $X\sim U(0,6\pi)$，求 $Y=\cos X$ 的密度。

$Y\in[-1,1]$。对 $-1<y<1$，令 $\theta=\arccos y\in(0,\pi)$。区间 $(0,6\pi)$ 包含三个完整周期，每个周期上余弦先减后增，所以所有原像恰为六个互不相同的点
\[
h_{k,1}(y)=\arccos y+2k\pi,\qquad
h_{k,2}(y)=2\pi-\arccos y+2k\pi,
\quad k=0,1,2.
\]
六个反函数分支的导数绝对值均为 $1/\sqrt{1-y^2}$，原密度均为 $1/(6\pi)$，于是
\begin{align*}
f_Y(y)
&=\sum_{k=0}^2
\left[
\frac1{6\pi\sqrt{1-y^2}}+
\frac1{6\pi\sqrt{1-y^2}}
\right]\\
&=\frac1{\pi\sqrt{1-y^2}},\qquad -1<y<1.
\end{align*}
在 $y=\pm1$ 处，原像为有限个点，概率为零；分支导数公式在端点不适用，可将密度另取为 $0$。所以一个完整版本为
\[
f_Y(y)=
\begin{cases}
\dfrac1{\pi\sqrt{1-y^2}},&-1<y<1,\\
0,&\text{其他}.
\end{cases}
\]
虽然密度在两端附近无界，但积分收敛，并且
\[
\int_{-1}^1\frac{dy}{\pi\sqrt{1-y^2}}
=\frac1\pi\left[\arcsin y\right]_{-1}^1=1.
\]

\subsection{例题：倒数的小数部分与可数个反函数分支}\label{note:c2:e25}
\basis{原笔记第2.4节，例“倒数的小数部分与可数个反函数分支”；源码\texttt{chapters/chapter2.tex}第2084行。}
设 $X\sim U(0,1)$，令 $Y$ 为 $1/X$ 的小数部分。记 $\lfloor u\rfloor$ 为不超过 $u$ 的最大整数，即向下取整，因而
\[
Y=\frac1X-\left\lfloor\frac1X\right\rfloor\in[0,1).
\]
在零概率事件 $\{X\notin(0,1)\}$ 上，可另定义 $Y=0$。

对 $0<y<1$，若 $Y=y$，则唯一的整数部分 $k=\lfloor1/X\rfloor$ 满足 $k\geq1$，且
\[
\frac1X=k+y
\quad\Longleftrightarrow\quad X=h_k(y)=\frac1{k+y}.
\]
反过来，每个 $k\geq1$ 给出的 $h_k(y)\in(0,1)$ 都是有效原像。各分支对应互不相交的区间
\[
I_k=\left(\frac1{k+1},\frac1k\right).
\]
这些区间只遗漏可数个点 $1/k$，其总概率为零，满足分支公式的全概率条件。由于
\[
|h_k'(y)|=\frac1{(k+y)^2},
\]
得到一个密度版本
\[
f_Y(y)=
\begin{cases}
\displaystyle\sum_{k=1}^{\infty}\frac1{(k+y)^2},&0\leq y<1,\\
0,&\text{其他}.
\end{cases}
\]
这里在 $y=0$ 处的值是另选的密度值；该点的原像是可数集，$P(Y=0)=0$。对 $0\leq y<1$，各项不超过 $1/k^2$，所以级数收敛。

利用非负求和与积分交换，再作望远镜求和，核验
\begin{align*}
\int_0^1 f_Y(y)\,dy
&=\sum_{k=1}^{\infty}\int_0^1\frac{dy}{(k+y)^2}\\
&=\sum_{k=1}^{\infty}\left(\frac1k-\frac1{k+1}\right)
=1.
\end{align*}
